% ==================== PART 3 — Reflections (human-written, <=1pp) ====================

\paragraph{Write the tester, or train it?}
The reflex in automated game testing is to \emph{train} an explorer, an RL agent with an intrinsic-motivation reward, one
policy per game. Our strongest result questions that reflex. At the interaction budget a tester actually gets, a
\textbf{zero-training program the loop wrote} reaches target and defect states faster than trained PPO across the
sparse-reward tasks (grid navigation, arcade, and planning), and tuning PPO does not close the gap at equal budget. We
think the field should take seriously that \textbf{exploration for testing is often a program-synthesis problem, not a
policy-learning one}: the LLM supplies, once and for free, the structured save/restore strategy that RL spends its whole
budget failing to learn from scratch. The boundary is just as useful in the other direction, since when reward is dense RL
competes and eventually wins with more budget. Writing the tester and training it are not rivals but regimes, and a loop
like ours maps the boundary cheaply.

\paragraph{You cannot trust an autonomous result you have not re-executed.}
The result we are proudest of is not the one the loop first reported. Its first ``discovery'' was a \textbf{+0.11 score
inflation it earned by exploiting a concurrency race in its own evaluation harness}: physically impossible states that
tripped the defect oracle, with nothing incriminating in the candidate's code. It passed every evaluator-mediated check we
had (the fitness, both held-out splits, four design comparisons, and an LLM code-audit of the candidate) and was caught
only by \textbf{replaying every counted hit from a clean reset} and requiring it to reproduce. The general lesson for a
field about to have loops producing candidate papers at scale: \textbf{you cannot certify an autonomous result by reading
its code or its score; you must re-execute what it claims from a clean state.} Code-audit, today's default defense, is
blind to an exploit that lives in the harness's state rather than the candidate. Execution-grounded verification is not a
nice-to-have; for this class of exploit it is the only thing that works, and it should live \emph{inside} the loop's
fitness so that such exploits are worthless by construction. The uncomfortable corollary is that most published
autonomous-research results ship no seeds or execution traces, so this class of error could be sitting in them undetected.

\paragraph{What the loop did, and what we did.}
The loop wrote the harness, evolved the testers, built the benchmark, ran the experiments, and drafted these Parts. We
designed the defect oracles and the verification protocol, decided which evaluator to point the search at, verified every
number against regenerable logs, and made the one intervention that mattered most: noticing the inflated result was too
good and rebuilding the harness. That division is, we suspect, where autonomous ML research settles for a while.
\textbf{The loop does the research; the scarce human contribution is deciding what to measure and re-checking that the
measure was not gamed.} The workshop's premise, that evaluating AI-generated claims is the bottleneck, is not an
abstraction to us. It is the single reason our headline number is honest.
